\documentclass[14pt, a4paper]{extarticle}
\usepackage[T2A]{fontenc}
\usepackage[utf8]{inputenc}
\usepackage[english, russian]{babel}
\usepackage{tempora}
\usepackage{longtable}
\usepackage{array}
\usepackage{ltablex}
\usepackage{ragged2e}
\usepackage{graphicx}
\usepackage{calc}
\usepackage[left=30mm, right=15mm, top=20mm, bottom=20mm, footskip=12mm]{geometry}

\setlength{\parindent}{0pt}
\setlength{\parskip}{6pt plus 2pt minus 1pt}

\newcolumntype{P}[1]{>{\RaggedRight\arraybackslash}p{#1 - 2\tabcolsep - 2\arrayrulewidth}}
\newcolumntype{C}[1]{>{\centering\arraybackslash}p{#1 - 2\tabcolsep - 2\arrayrulewidth}}


\begin{document}

\pagestyle{empty}

% ------------------------- Заголовок документа -------------------------

\begin{center}
    {\small МИНИСТЕРСТВО НАУКИ И ВЫСШЕГО ОБРАЗОВАНИЯ РОССИЙСКОЙ ФЕДЕРАЦИИ}

    \resizebox{\textwidth}{!}{\fontsize{8pt}{10pt}\selectfont ФЕДЕРАЛЬНОЕ
        ГОСУДАРСТВЕННОЕ АВТОНОМНОЕ ОБРАЗОВАТЕЛЬНОЕ УЧРЕЖДЕНИЕ ВЫСШЕГО ОБРАЗОВАНИЯ}

    {\small \textbf{«МОСКОВСКИЙ ПОЛИТЕХНИЧЕСКИЙ УНИВЕРСИТЕТ»}} \\
    {\small \textbf{(МОСКОВСКИЙ ПОЛИТЕХ)}}
\end{center}

\begin{tabularx}{\textwidth}{@{} p{0.53\textwidth} p{0.43\textwidth} @{}}
     &
    \raggedright
    УТВЕРЖДАЮ                          \\
    заведующая кафедрой                \\
    «Инфокогнитивные технологии»       \\[0.3cm]

    \rule{3cm}{0.4pt} / Е.~А.~Пухова / \\[0.1cm]

    «\rule{0.8cm}{0.4pt}» \rule{3cm}{0.4pt} 2026~г.
\end{tabularx}

% -------------------------- Общая информация --------------------------

\begin{center}

    { \large \textbf{ЗАДАНИЕ}}

    \textbf{на выполнение курсовой работы (проекта)}

    Дудченко Алёне Алексеевне, % нужно указать в дательном падеже (задание кому?)

\end{center}

обучающемуся группы 241-326,

направления подготовки 09.03.01 «Информатика и вычислительная техника»

по дисциплине «Разработка веб-приложений»

на тему «Разработка веб-приложения для оптимизации логистических
маршрутов службы экспресс-доставки»

1. Исходные данные к работе (проекту): информационные ресурсы в сети
интернет, научные публикации в открытой печати.

2. Содержание задания по курсовой работе (проекту) -- перечень вопросов,
подлежащих разработке:

% ------------------------------- Задачи -------------------------------

{
\renewcommand{\arraystretch}{1.2}
\small
\begin{longtable}{|P{0.5\textwidth}|C{0.11\textwidth}|C{0.15\textwidth}|C{0.24\textwidth}|}

    \hline
    \textbf{Разрабатываемый вопрос}                                                                  & \textbf{Объем, \%} & \textbf{Срок} & \textbf{Примечание} \\
    \hline
    \endhead

    \hline
    \endfoot % Этот блок сработает при разрыве (внизу страницы)

 \multicolumn{4}{|P{\textwidth}|}{\textbf{Раздел 1. Анализ предметной области}} \\ \hline
    Задача 1.1. Сравнительный анализ существующих платформ для поиска учебных партнеров и выявление их недостатков & 10 & 30.03.2026 & \\ \hline
    Задача 1.2. Обоснование выбора стека технологий (Python, Flask, PostgreSQL) & 7 & 06.04.2026 & \\ \hline
    Задача 1.3. Формулировка целей, задач и актуальности проекта & 8 & 13.04.2026 & \\ \hline

    \multicolumn{4}{|P{\textwidth}|}{\textbf{Раздел 2. Проектирование веб-приложения}} \\ \hline
    Задача 2.1. Анализ целевой аудитории и проектирование сценариев взаимодействия (User Stories) & 5 & 20.04.2026 & \\ \hline
    Задача 2.2. Разработка диаграммы вариантов использования (UML Use Case) & 6 & 27.04.2026 & \\ \hline
    Задача 2.3. Проектирование модели данных (ER-диаграмма) для системы из 4-х сущностей & 7 & 04.05.2026 & \\ \hline
    Задача 2.4. Разработка макетов интерфейса (Wireframes) основных страниц сервиса & 5 & 11.05.2026 & \\ \hline

    \multicolumn{4}{|P{\textwidth}|}{\textbf{Раздел 3. Разработка веб-приложения}} \\ \hline
    Задача 3.1. Реализация базовой структуры и механизмов аутентификации и авторизации пользователей & 8 & 18.05.2026 & \\ \hline
    Задача 3.2. Разработка CRUD-интерфейса для управления профилями и реализации системы прав доступа по ролям & 9 & 25.05.2026 & \\ \hline
    Задача 3.3. Реализация алгоритма агрегации данных и подбора партнеров на основе профиля компетенций & 9 & 01.06.2026 & \\ \hline
    Задача 3.4. Реализация модуля загрузки и выгрузки дополнительных учебных материалов & 8 & 08.06.2026 & \\ \hline
    Задача 3.5. Создание пользовательского интерфейса с использованием шаблонизатора Jinja2 & 8 & 10.06.2026 & \\ \hline

    \multicolumn{4}{|P{\textwidth}|}{\textbf{Раздел 4. Оформление итогов работы}} \\ \hline
    Задача 4.1. Тестирование функций системы и управление версиями в репозитории Git & 3 & 11.06.2026 & \\ \hline
    Задача 4.2. Развертывание приложения на хостинге & 3 & 12.06.2026 & \\ \hline
    Задача 4.3. Оформление пояснительной записки согласно требованиям в формате LaTeX & 4 & 13.06.2026 & \\ \hline
\end{longtable}
}

% -------------- Сведения о выдаче и принятии задания -------------------

Руководитель курсовой работы (проекта): \\ старший преподаватель кафедры
«Инфокогнитивные технологии»

{
\renewcommand{\arraystretch}{2}
\noindent
\begin{tabularx}{\textwidth}{@{} P{0.45\textwidth} P{0.3\textwidth} P{0.25\textwidth} @{}}

    «\rule{0.8cm}{0.4pt}» \quad \rule{2cm}{0.4pt} \quad 2026~г. &
    \centering \rule{3.5cm}{0.4pt} \par \footnotesize (подпись) &
    \hfill А.~С.~Кружалов                                         \\
\end{tabularx}

\begin{tabularx}{\textwidth}{@{} P{0.45\textwidth} P{0.55\textwidth} @{}}
    Дата выдачи задания           & 
    \hfill «\underline{31}» \quad \underline{марта} \quad 2026~г. \\

    Дата сдачи выполненной работы & 
    \hfill «\rule{0.8cm}{0.4pt}» \quad \rule{2cm}{0.4pt} \quad 2026~г. \\
\end{tabularx}
}

Задание принял к исполнению

{
\renewcommand{\arraystretch}{1.5}
\begin{tabularx}{\textwidth}{@{} P{0.45\textwidth} P{0.3\textwidth} P{0.25\textwidth} @{}}

    «\rule{0.8cm}{0.4pt}» \quad \rule{2cm}{0.4pt} \quad 2026~г. &
    \centering \rule{3.5cm}{0.4pt} \par \footnotesize (подпись) &
    \hfill А. А. Дудченко                                           \\
\end{tabularx}
}

\end{document}
